\documentclass[11pt,a4paper]{article}
\usepackage[utf8]{inputenc}
\usepackage[T1]{fontenc}
\usepackage[english]{babel}
\usepackage{geometry}
\usepackage{graphicx}
\usepackage{hyperref}
\usepackage{listings}
\usepackage{xcolor}
\usepackage{booktabs}
\usepackage{amsmath}
\usepackage{fancyhdr}
\usepackage{titlesec}
\usepackage{enumitem}


\geometry{margin=2.5cm}

\definecolor{codeblue}{RGB}{0,102,204}
\definecolor{codegray}{RGB}{128,128,128}
\definecolor{codegreen}{RGB}{0,140,0}
\definecolor{codered}{RGB}{180,0,0}
\definecolor{codebg}{RGB}{248,248,248}

\lstset{
  backgroundcolor=\color{codebg},
  basicstyle=\ttfamily\small,
  keywordstyle=\color{codeblue}\bfseries,
  commentstyle=\color{codegreen}\itshape,
  stringstyle=\color{codered},
  numberstyle=\tiny\color{codegray},
  numbers=left,
  stepnumber=1,
  frame=single,
  rulecolor=\color{codegray},
  breaklines=true,
  captionpos=b,
  tabsize=4,
  showstringspaces=false
}

\lstdefinelanguage{Lua}{
  keywords={and,break,do,else,elseif,end,false,for,function,if,in,
            local,nil,not,or,repeat,return,then,true,until,while,
            pcall,print,tostring},
  morestring=[b]",
  morestring=[b]',
  morecomment=[l]{--},
  morecomment=[s]{--[[}{]]},
  sensitive=true
}

\pagestyle{fancy}
\fancyhf{}
\lhead{\textit{IoT Seismograph — Report v1}}
\rhead{\textit{O. Denis}}
\cfoot{\thepage}

\hypersetup{
  colorlinks=true,
  linkcolor=codeblue,
  urlcolor=codeblue,
  citecolor=codeblue
}

\title{
  \vspace{-1cm}
  \textbf{\Large IoT Embedded Alert System} \\[0.4em]
  \large Proof of Concept — Report \\[0.2em]
  \normalsize ESP32 + Hall Sensor + Telegram Alert
}
\author{O. Denis \\ \small [University] — IoT Training Program}
\date{2026}

\begin{document}

\maketitle
\thispagestyle{fancy}
\tableofcontents
\vspace{1em}
\hrule
\vspace{1em}

% -------------------------------------------------------------------
\section{Project Identification}

\subsection{Solution Name}
\textbf{IoT Seismograph} — Real-time seismic detection system using a Hall-effect sensor and an oscillating magnet mounted on a mechanical pendulum, with instant alerting via Telegram Bot API.

\subsection{Group}
\begin{itemize}[noitemsep]
  \item O. Denis
\end{itemize}

\subsection{Proposed Applications}
\begin{enumerate}[noitemsep]
  \item \textbf{Intelligent Seismic Detector (ISD)} — critical vibration monitoring
  \item Door Opening Detector — building security / access audit
  \item Vehicle Counter for Parking — smart access automation
\end{enumerate}

% -------------------------------------------------------------------
\section{Context and Problem Statement}

\subsection{Use Context}
Target deployment: environments where undetected vibrations can cause hardware damage before any human intervention. Primary scenario: a data center where undetected micro-vibrations cause intermittent hard drive failures. Administrators only discover the issue after data loss. An early warning system enables intervention before hardware failure.

Secondary scenarios:
\begin{itemize}[noitemsep]
  \item Hospital or laboratory (sensitive equipment protection)
  \item Underground parking (heavy vehicle vibration detection)
  \item Electrical substation
\end{itemize}

\subsection{General Objective}
Demonstrate the feasibility of an autonomous, real-time, low-cost IoT system capable of detecting seismic events (P-waves ideally) with instant notification via a secure channel (Telegram Bot API).

% -------------------------------------------------------------------
\section{Solution Value}

\subsection{Key Benefits}
\begin{itemize}[noitemsep]
  \item \textbf{Automatic detection}: zero manual intervention, 24/7 operation
  \item \textbf{Real-time alert}: notification within 5 s via Telegram (WiFi)
  \item \textbf{Low cost}: commodity hardware (ESP32 + KY-003 Hall sensor)
  \item \textbf{Noise discrimination}: threshold-based filter separates genuine seismic patterns from electrical rebounds and industrial machines
  \item \textbf{Resilience}: local serial output continues if WiFi is unavailable
\end{itemize}

\subsection{Alert Use Cases}
\begin{enumerate}[noitemsep]
  \item \textit{Major seismic alert} → urgent Telegram message, immediate response required
  \item \textit{Moderate pre-alert} → informational notification, manual investigation
  \item \textit{Vibration anomaly} → distinguishes isolated shock (noise) from sustained seismic event (duration $>$ 3 s)
\end{enumerate}

% -------------------------------------------------------------------
\section{Hardware and Software Resources}

\subsection{Hardware Components}

\begin{table}[h!]
\centering
\begin{tabular}{ll}
\toprule
\textbf{Component} & \textbf{Role} \\
\midrule
ESP32 & Main MCU, WiFi, GPIO \\
KY-003 Hall Magnetic Field Sensor & Oscillation detector \\
Oscillating magnet (pendulum) & Mechanical seismic transducer \\
128×64 OLED screen & Local status display \\
Spring / suspension wire & Mechanical pendulum assembly \\
External battery / USB power & Power supply \\
Enclosure / support & Physical mounting \\
\bottomrule
\end{tabular}
\caption{Hardware bill of materials}
\end{table}

\subsection{Software Environment}
\begin{itemize}[noitemsep]
  \item \textbf{Language}: Lua 5.3.4 on LuaRTOS
  \item \textbf{Modules}: \texttt{pio}, \texttt{net.wf}, \texttt{net.curl}, \texttt{tmr}, \texttt{thread}, \texttt{os}, \texttt{string}
  \item \textbf{External API}: Telegram Bot API — \texttt{https://api.telegram.org/bot[TOKEN]/sendMessage}
\end{itemize}

% -------------------------------------------------------------------
\section{Technical Details of the PoC}

\subsection{System Architecture}

The system operates in two layers:

\begin{enumerate}[noitemsep]
  \item \textbf{Sensor layer} (30--50 ms polling): reads the Hall pin, counts falling edges (magnet passes), evaluates the count every 3--7 seconds against thresholds.
  \item \textbf{Communication layer} (async thread): picks up alert level from shared variable, constructs and sends HTTPS GET request to Telegram API, enforces 30-second anti-spam cooldown.
\end{enumerate}

\subsection{Detection Algorithm}

\begin{lstlisting}[language=Lua, caption={Sensor thread core logic}]
-- Analysis every 3 seconds
if os.time() - start_time >= 3 then
    if detections >= 5 and not send_in_progress then
        local level = "LOW"
        if detections >= 10 then level = "MODERATE" end
        if detections >= 20 then level = "HIGH"     end

        -- Pass alert to Telegram thread
        alert_level        = level
        alert_oscillations = detections
    end
    detections = 0
    start_time = os.time()
end
\end{lstlisting}

\subsection{Alert Thresholds}

\begin{table}[h!]
\centering
\begin{tabular}{lll}
\toprule
\textbf{Parameter} & \textbf{Value} & \textbf{Justification} \\
\midrule
Min threshold $S_{min}$ & 5 oscillations / 3 s & Below = electrical rebound noise \\
Max threshold $S_{max}$ & 40 oscillations / s & Above = industrial fans (50--60 Hz) \\
Confirmation duration & 3 s & Isolates single shock vs sustained seism \\
Analysis cycle & 7 s (debug) / 3 s (prod) & Balances latency and CPU load \\
Anti-spam cooldown & 30 s minimum & Prevents 10+ messages for one event \\
\bottomrule
\end{tabular}
\caption{Detection parameter justification}
\end{table}

\subsection{Noise vs Real Event Discrimination}

\begin{itemize}[noitemsep]
  \item \textbf{Object falling on table} (noise): peak of 50 fronts/s for 0.2 s → filtered ($> S_{max}$), then immediate silence
  \item \textbf{Seismic event} (real): 8--15 fronts/s for 5--10 s continuously → confirmed, chaotic pattern
  \item \textbf{Industrial motor at 50 Hz} (noise): constant 50 fronts/s for hours → filtered ($> S_{max}$)
\end{itemize}

\subsection{Security Analysis}

The Telegram alert travels in plaintext over WiFi. Identified risks and mitigations:

\begin{table}[h!]
\centering
\begin{tabular}{p{3.5cm}p{4cm}p{5.5cm}}
\toprule
\textbf{Threat} & \textbf{Risk} & \textbf{Mitigation} \\
\midrule
Replay attack & Attacker resends captured message & Sequence number counter (16-bit) \\
Message forgery & Fake alert injected & HMAC-SHA256 signature (requires porting crypto lib) \\
Token exposure & Firmware extraction → bot abuse & Store in encrypted EEPROM; revoke after demo \\
WPA2 interception & Token visible in cleartext & Enforce WPA3; add VPN (heavy for ESP32) \\
\bottomrule
\end{tabular}
\caption{Security threat matrix}
\end{table}

Sequence number anti-replay (Lua pseudocode):
\begin{lstlisting}[language=Lua, caption={Anti-replay sequence counter}]
local sequence = 0
function send_alert(level, count)
    sequence = (sequence + 1) % 65536
    local msg = level .. "|" .. count .. "|" .. sequence
    -- Server-side: reject if sequence <= last_seen_sequence
end
\end{lstlisting}

% -------------------------------------------------------------------
\section{Limitations}

\subsection{Physical Limitations}
\begin{itemize}[noitemsep]
  \item \textbf{Single-axis}: Hall sensor detects only vertical movement. Lateral tremors (East-West, North-South) are missed — 66\% sensitivity reduction for horizontal seismic waves. Mitigation: add 2 additional Hall sensors on perpendicular axes.
  \item \textbf{No amplitude measurement}: binary detection (0/1). Richter scale approximated to 3 levels only.
\end{itemize}

\subsection{Network Limitations}
\begin{itemize}[noitemsep]
  \item \textbf{WiFi dependency}: no Telegram notification if WiFi goes down.
  \item \textbf{Telegram latency}: HTTP REST adds 2--5 s. Acceptable since seismic events last 10--30 s.
\end{itemize}

\subsection{Operational Limitations}
\begin{itemize}[noitemsep]
  \item \textbf{No distributed consensus}: a single node vibrating triggers a global alert. Door slam in one room → all admins notified. Mitigation: require $N$ of $M$ nodes to agree before sending.
  \item \textbf{Rate limiting}: 10+ nodes sending simultaneously risk hitting Telegram's 30 msg/s global limit. Mitigation: aggregate into one message (\texttt{Node1:OK|Node2:ALERT|Node3:OK}).
\end{itemize}

% -------------------------------------------------------------------
\section{Conclusion}

This PoC validates the core feasibility of an ESP32-based seismic monitoring system at minimal cost. The dual-thread architecture successfully decouples latency-sensitive sensor polling from blocking network I/O. The threshold-based classifier provides acceptable noise discrimination for a data-center deployment scenario. The primary gap remains mechanical: a single-axis, frequency-only sensor cannot produce a calibrated Richter magnitude. This would require a tri-axis MEMS accelerometer (e.g., MPU-6050) as a hardware upgrade path.

\end{document}
